\section{Experimental Setup}
\label{sec:experiments}

\paragraph{Training Data.}
We sample a 132M-token subset of FineWeb~\citep{penedo_2024_fineweb} as the real-data source.
To simulate web-scale synthetic contamination, we generate continuations with Qwen3-0.6B~\citep{yang_2025_qwen3}---the same architecture used for continued pretraining---and inject them at ratio $\rho = 0.4$ via uniform random mixing.
This self-contamination design creates a challenging detection scenario where generator and trainee share the same distribution.

\paragraph{Detector.}
We train a TF-IDF + Logistic Regression classifier on a held-out labeled split, achieving AUC\,=\,0.939 with calibrated ECE\,=\,0.048 after temperature scaling ($T = 0.204$).
The resulting $q$-scores have mean 0.60, standard deviation 0.37, and median 0.71, with 46.3\% of samples in the uncertain range $[0.2, 0.8]$.
We use this lightweight detector rather than neural alternatives (e.g., ModernBERT at AUC\,${\geq}$\,0.98 in \citealt{drayson_2025_detection_collapse}) to test whether score \emph{calibration} matters more than raw discrimination for reweighting.

\paragraph{Training.}
We continue pretraining Qwen3-0.6B from its pretrained weights for 6{,}103 steps (batch size 4, gradient accumulation 4, effective batch 16) with AdamW, cosine schedule (peak lr\,=\,$5{\times}10^{-5}$), max sequence length 2048, and bf16 precision.
All main comparisons use three random seeds (42, 123, 456).

\paragraph{Evaluation.}
We report held-out perplexity (PPL) on a held-out subset of FineWeb (5{,}000 documents) containing only human-authored text, via sliding-window evaluation (stride 1024).
All evaluations are post-hoc with a unified pipeline; training involves no inline evaluation.

\paragraph{Methods.}
We compare five approaches:
\begin{itemize}[nosep,leftmargin=*]
  \item \textbf{Uniform}: $w_j = 1$ for all $j$ (no contamination handling).
  \item \textbf{Grid-$b$}: fixed $b \in \{0.5, 1.0, 1.5, 2.0, 3.0, 5.0\}$, best selected post-hoc by held-out PPL.
        Represents validation-based selection~\citep{drayson_2025_detection_collapse}.
  \item \textbf{ESS-only}: $b^* = \max\{b : \mathrm{ESS}(b) \geq n_{\min}\}$, the maximum-intervention variant retaining only the stability constraint, ablating $\alpha_{\mathrm{eff}}$.
  \item \textbf{Golden Ratio}: $w = 1/\varphi \approx 0.618$ for synthetic ($q_j > 0.5$), $w = 1$ otherwise~\citep{he_2025_golden_ratio}.
        Requires known provenance labels (oracle).
  \item \textbf{Adaptive} (ours): per-batch $b^*$ from the dual-constraint formulation (\S\ref{sec:method}), with $\tau = 0.8$ and $n_{\min} = 2$ (half the micro-batch size).
\end{itemize}
